% ============================================================
% PROPOSED METHODOLOGY
% ============================================================
\section{Proposed Methodology}

\subsection{System Overview}

The system has two halves (Figure~\ref{fig:arch}). The compiler half is a conventional pipeline (scanner, parser,
semantic analyser, intermediate-code generator, optimiser) for a language called MedScript, and every phase reports
diagnostics with a stable code, a position and a message. The learning half condenses the tokens, syntax tree and
diagnostics into a fixed-length feature vector, and a classifier maps it to one of ten classes: \emph{valid} or one of
nine error families. A conditional GAN trained on the labelled vectors supplies extra records of rare classes at
training time only. What gets deployed is the compiler plus a random forest.

\begin{figure}[!htbp]
\centering
\resizebox{0.92\textwidth}{!}{%
\begin{tikzpicture}[
  box/.style={rectangle,rounded corners=3pt,draw,minimum height=0.9cm,minimum width=2.1cm,align=center,font=\small,fill=blue!6},
  ml/.style={rectangle,rounded corners=3pt,draw,minimum height=0.9cm,minimum width=2.4cm,align=center,font=\small,fill=orange!12},
  gan/.style={rectangle,rounded corners=3pt,draw,minimum height=0.9cm,minimum width=2.4cm,align=center,font=\small,fill=red!8},
  arr/.style={-{Stealth[length=2.2mm]},thick},
  dsh/.style={-{Stealth[length=2mm]},dashed}
]
\node[box] (src) at (0,0) {MedScript\\source};
\node[box] (lex) at (2.9,0) {Lexer\\(scanner)};
\node[box] (par) at (5.8,0) {Parser\\(recursive descent)};
\node[box] (sem) at (8.7,0) {Semantic\\analyser};
\node[box] (ir)  at (11.6,0) {IR generator\\+ optimiser};
\draw[arr] (src) -- (lex);
\draw[arr] (lex) -- (par);
\draw[arr] (par) -- (sem);
\draw[arr] (sem) -- (ir);

\node[box,fill=green!8,minimum width=4.2cm] (diag) at (5.8,-2.3) {Diagnostics + features\\(101-dimensional vector)};
\draw[dsh] (lex.south) -- ++(0,-0.6) -| ($(diag.north)+(-1.0,0)$);
\draw[dsh] (par.south) -- (diag.north);
\draw[dsh] (sem.south) -- ++(0,-0.6) -| ($(diag.north)+(1.0,0)$);

\node[ml] (clf) at (10.0,-2.3) {Error-family\\classifier};
\draw[arr] (diag) -- (clf);
\node[box,fill=gray!10] (rep) at (13.0,-2.3) {Report:\\diagnostics +\\root cause};
\draw[arr] (clf) -- (rep);
\draw[arr] (ir) -- (13.0,0) -- (rep.north);

\node[box,fill=gray!10] (gen) at (0.8,-4.6) {Program generator\\+ 46 fault operators};
\node[ml,fill=orange!6] (corp) at (4.2,-4.6) {Labelled feature\\vectors};
\node[gan] (cg) at (7.4,-4.6) {Conditional\\GAN};
\node[gan,fill=red!4] (syn) at (10.4,-4.6) {Synthetic\\vectors};
\draw[arr] (gen) -- (corp);
\draw[arr] (corp) -- (cg);
\draw[arr] (cg) -- (syn);
\draw[arr] (syn) -- (clf.south);
\draw[dsh] (corp.north) -- ++(0,0.9) -| (diag.south west);
\node[font=\scriptsize,align=center] at (4.2,-3.35) {same extractor};
\end{tikzpicture}}
\caption{System architecture. Solid arrows are data flow, dashed arrows are diagnostics and features collected from
the compiler phases. The GAN branch (bottom) is only used to enlarge the training set.}
\label{fig:arch}
\end{figure}

\subsection{The MedScript Language}

MedScript is small on purpose: enough to write the programs a clinical-informatics student would write, and enough
machinery to make the domain checks meaningful.

\begin{itemize}
    \item \textbf{Units are part of the type.} A numeric type is \code{int}, \code{float}, or a float with a unit such
    as \code{kg}, \code{mg/kg} or \code{kg/metre\textasciicircum 2}. Literals carry units: \code{38.0 degC}. The
    vocabulary has eleven units (\code{kg}, \code{lb}, \code{gram}, \code{mg}, \code{ml}, \code{hour}, \code{bpm},
    \code{mmHg}, \code{degC}, \code{degF}, \code{metre}), each with a plausible literal range such as 20 to 45 for \code{degC}.
    \item \textbf{Inputs carry metadata.} \code{input weight: kg range(2, 300);} declares a read-only patient value with
    a physiological range, and \code{phi} marks protected health information.
    \item \textbf{Optional values.} A trailing \code{?} means the value may be missing and must be tested with
    \code{present(x)} before use.
    \item \textbf{Two sinks}, \code{print} and \code{export}, which reject PHI unless it went through \code{anonymize}.
    \item \textbf{Ordinary control flow}: \code{if}/\code{else}, \code{while}, typed functions and \code{let} bindings.
\end{itemize}

Listing~\ref{lst:ok} shows a program the compiler accepts: a capped weight-based dose, a BMI, a fever alert, a sepsis
screen, a guarded check on a possibly missing lab value, and an anonymised export.

\begin{lstlisting}[language=MedScript,caption={A valid MedScript program (\texttt{examples/triage\_ok.med}).},label={lst:ok}]
input patient_name: string phi;
input weight: kg range(2, 300);
input temp: degC range(30, 43);
input hr: bpm range(20, 250);
input lactate: mg? range(0, 400);
const MAX_DOSE: mg = 1000 mg;

func dose_for(w: kg, rate: mg/kg) -> mg {
    let d: mg = rate * w;
    if (d > MAX_DOSE) { return MAX_DOSE; }
    return d;
}
let dose: mg = dose_for(weight, 15 mg/kg);
let fever = false;
if (temp > 38.0 degC) { fever = true; print("fever alert"); }
if (hr > 100 bpm && fever) { print("sepsis screen"); }
if (present(lactate)) {
    if (lactate > 200 mg) { print("high lactate"); }
}
export(anonymize(patient_name));
print(dose);
\end{lstlisting}

\subsection{Lexical and Syntax Analysis}

The scanner is a hand-written, single-pass, longest-match recogniser for identifiers, 16 keywords, 11 unit words,
numbers, strings, operators and \code{//} comments. A numeric lexeme is collected as the longest run of letters, digits
and dots and only then checked against $[0\text{-}9]^+(.[0\text{-}9]^+)?$, so \code{3.4.5} is reported as one bad
number instead of a number plus junk. The three lexical errors (illegal character, malformed number, unterminated
string; \code{LEX001} to \code{LEX003}) are all recoverable, so later phases still run. That matters for
classification: a lexical error does not hide semantic faults, but a syntax error does.

The grammar is LL(1) except in two places. A statement starting with an identifier needs one extra token to choose
between assignment and call, and after a number and a unit, a \code{/} belongs to the unit only if another unit follows.
Expressions use precedence climbing over six levels. Error recovery combines two mechanisms. A missing \code{;},
\code{)} or \code{\}} is reported as \code{SYN001} at the end of the previous token, where a programmer expects to see
it, and parsing continues as if the token were there. When an expression cannot start (\code{SYN002}), the parser skips
to the next semicolon or statement keyword. At most one diagnostic is issued per position, which removes most duplicate
cascades. If any syntax error was reported, semantic analysis is skipped.

\subsection{Semantic Analysis}

The analyser walks the tree once with a stack of scopes and a small flow state. This is where the clinical knowledge
lives. Functions are collected in a first pass, so calls can come before definitions, and the usual name errors
(undeclared, redeclared, uninitialised, assignment to an input) are reported.

\paragraph{Types and units.} A numeric type is a pair $\tau[d]$ with $\tau\in\{\texttt{int},\texttt{float}\}$ and
$d\in\mathbb{Z}^{U}$ a dimension vector holding one exponent per base unit. For example \code{mg/kg} has $+1$ for
\code{mg} and $-1$ for \code{kg}. The arithmetic rules are
\begin{equation}
\frac{e_1:\tau_1[d_1]\quad e_2:\tau_2[d_2]\quad d_1=d_2}{e_1\pm e_2:\tau[d_1]}\qquad
\frac{e_1:\tau_1[d_1]\quad e_2:\tau_2[d_2]}{e_1\cdot e_2:\tau[d_1+d_2]}\qquad
\frac{e_1:\tau_1[d_1]\quad e_2:\tau_2[d_2]}{e_1/e_2:\texttt{float}[d_1-d_2]}
\end{equation}
and comparison needs $d_1=d_2$. So \code{rate * w} with \code{rate}:\code{mg/kg} and \code{w}:\code{kg} comes out as
\code{mg}, which is what \code{let d: mg} expects, while \code{weight + 5 lb} fails because \code{kg} and \code{lb} are
distinct and there is no implicit conversion. A pure dimension mismatch is \code{SEM005}; a base-type, arity or
condition mismatch is \code{SEM004}, \code{SEM006} or \code{SEM008}. Keeping them apart lets the classifier tell the
\emph{type} and \emph{unit} families apart later.

\paragraph{Ranges.} Every numeric value can carry an interval $[l,h]$ from an input's declared range or from interval
arithmetic on constants. A unit literal outside its plausible range (\code{380 degC}) is flagged, and so is a comparison
of a ranged input against a constant outside $[l,h]$, since it is then always true or always false. Both are
\code{DOM001}. A constant zero divisor is \code{DOM004}; a divisor whose interval contains zero is \code{DOM005} unless
a branch condition such as \code{x > 0} has established it is non-zero.

\paragraph{PHI taint.} Let $\mathrm{t}(e)\in\{0,1\}$ say whether $e$ may hold protected health information. Then
$\mathrm{t}(x)=\mathit{phi}(x)$, $\mathrm{t}(e_1\odot e_2)=\mathrm{t}(e_1)\lor\mathrm{t}(e_2)$, a function call is
tainted if any argument is, and $\mathrm{t}(\texttt{anonymize}(e))=0$. A tainted argument to a sink is \code{DOM002}.
The analysis is conservative: it never misses a leak, but a function that throws away its tainted argument still counts
as propagating it.

\paragraph{Optional values and flow.} Using an optional variable is \code{DOM003} unless a presence fact holds. Facts
come from \code{present(x)} guards, from \code{present(x) \&\& c}, from the else branch of \code{!present(x)}, and from
early exits like \code{if (!present(x)) \{ return 0; \}}. At a join the facts of both branches are intersected. The
same state tracks definite assignment, functions that can fall off the end without returning (\code{SEM007}),
unreachable statements (\code{W001}) and constant conditions (\code{W003}).

The analyser has 23 diagnostic codes, and each belongs to exactly one error family. These families are the label space
of the classifier (Table~\ref{tab:families}).

\begin{table}[!htbp]
\centering
\caption{The ten classes and their diagnostic codes.}
\label{tab:families}
\small
\begin{tabular}{clp{9.6cm}}
\toprule
\textbf{Id} & \textbf{Class} & \textbf{Diagnostic codes / meaning}\\
\midrule
0 & valid & no error (warning \code{W002} allowed)\\
1 & lexical & \code{LEX001} to \code{LEX003}: illegal character, malformed number, unterminated string\\
2 & syntax & \code{SYN001} to \code{SYN003}: missing token, unexpected token, unexpected end of file\\
3 & name & \code{SEM001}, \code{SEM002}, \code{SEM003}, \code{SEM009}: undeclared, redeclared, uninitialised, assignment to input\\
4 & type & \code{SEM004}, \code{SEM006}, \code{SEM008}: base-type, arity and condition errors\\
5 & unit & \code{SEM005}: unit mismatch\\
6 & range & \code{DOM001}: implausible value or threshold outside declared range\\
7 & phi & \code{DOM002}: PHI reaching a sink\\
8 & missing & \code{DOM003}: optional value used without a guard\\
9 & logic & \code{SEM007}, \code{DOM004}, \code{DOM005}, \code{W001}, \code{W003}, and rule errors the compiler cannot see\\
\bottomrule
\end{tabular}
\end{table}

\subsection{Intermediate Code and Optimisation}

An error-free tree is lowered to three-address code. Locals are renamed uniquely (\code{x@3}) so shadowing cannot
confuse later passes, short-circuit operators become conditional jumps, and calls become \code{param}/\code{call}
sequences. Six classic passes~\cite{muchnick1997,aho2006} run until the code stops shrinking (at most four rounds):
constant folding and propagation within basic blocks, algebraic simplification ($x+0$, $x\cdot 1$), local
common-subexpression elimination by value numbering, unreachable-code removal, dead-store elimination, and copy fusion
($t=a\,\mathit{op}\,b;\ v=t \Rightarrow v=a\,\mathit{op}\,b$). To check that the passes do not change behaviour, a small
virtual machine runs the original and optimised code on random inputs drawn from the declared ranges (with about 30\% of
optional inputs left missing) and requires identical \code{print} and \code{export} events.

\subsection{Feature Extraction and Labelling}

The classifier never sees source text. For each program the extractor builds a vector of \NFeat{} numbers in four
groups (Table~\ref{tab:features}). The diagnostic group carries the compiler's knowledge. The pattern group is there for
faults that produce no diagnostic: for each comparison, it checks whether nearby identifiers or strings contain a
``high-risk'' word (\emph{fever}, \emph{tachycardia}, \emph{overdose}, \ldots) or a ``low-risk'' word
(\emph{hypothermia}, \emph{bradycardia}, \ldots) and counts the four combinations of word type and operator direction. A
correct program mostly has ``high with $>$'' and ``low with $<$''; the other two point to a reversed rule.

\begin{table}[!htbp]
\centering
\caption{Feature groups.}
\label{tab:features}
\small
\begin{tabular}{lrp{10.2cm}}
\toprule
\textbf{Group} & \textbf{Dim.} & \textbf{Content}\\
\midrule
Token & 30 & log-counts of token classes, program length, brace and parenthesis balance, semicolons per statement\\
AST & 25 & log-counts of node types, tree size and depth, optional/PHI/ranged inputs, sinks, \code{anonymize} and \code{present} calls\\
Diagnostics & 40 & count of each of the 23 codes, error and warning totals, semantic-phase flag, family and phase of the first error, number of distinct families, worst-family rank\\
Pattern & 6 & comparison direction $\times$ alert vocabulary (four counts), number of alert rules, number of mismatched rules\\
\bottomrule
\end{tabular}
\end{table}

Each program has one label. With two injected faults, the label is the family with higher root-cause priority:
syntax $>$ lexical $>$ phi $>$ range $>$ unit $>$ missing $>$ type $>$ name $>$ logic. This order is a modelling choice.
Syntax comes first because it hides everything else, privacy and range next because they are the most safety-critical,
and logic last because it is the most mixed family.

\subsection{Conditional GAN}

Let $x\in\mathbb{R}^{101}$ be a standardised feature vector and $c$ the one-hot class. The generator maps noise
$z\sim\mathcal{N}(0,I_{32})$ and $c$ to $\tilde{x}=G(z,c)$, and the discriminator gives a score $D(x,c)$. Both are MLPs
with two hidden layers of 128 units and LeakyReLU(0.2), written in PyTorch~\cite{paszke2019}. We train two variants. The
BCE version plays the conditional minimax game
\begin{equation}
\min_G\max_D\ \mathbb{E}\big[\log D(x,c)\big]+\mathbb{E}\big[\log\big(1-D(G(z,c),c)\big)\big]
\end{equation}
with the non-saturating generator loss and real targets smoothed to 0.9~\cite{salimans2016}. The WGAN-GP critic
minimises
\begin{equation}
L_D=\mathbb{E}\big[D(\tilde{x},c)\big]-\mathbb{E}\big[D(x,c)\big]+\lambda\,\mathbb{E}_{\hat{x}}\Big[\big(\lVert\nabla_{\hat{x}}D(\hat{x},c)\rVert_2-1\big)^2\Big]
\end{equation}
with $\hat{x}$ sampled between real and fake points, $\lambda=10$ and five critic steps per generator step
\cite{gulrajani2017}. Both use Adam~\cite{kingma2015} and batch size 64 (2000 generator iterations for BCE, 1000 for
WGAN-GP). The GAN is trained on all classes jointly, on the training part of each experiment only. For augmentation,
every class below a target count is topped up with samples from $G(z,c)$, mapped back to the original scale and clipped
to the training range.

\subsection{Classifiers, Baselines and Anomaly Scoring}

We train four scikit-learn~\cite{pedregosa2011} classifiers: a 300-tree random forest~\cite{breiman2001}, an MLP (128
and 64 hidden units), logistic regression and histogram gradient boosting. Two rule baselines use only the diagnostics.
\emph{First error} predicts the family of the first reported error, which is what a developer reading top to bottom
would conclude. \emph{Most severe} predicts the highest-priority family among all errors. The second one is strong on
purpose, because it is the labelling rule itself, so anything the learned models gain over it comes from cases the rule
cannot express.

For anomaly detection a single-class GAN is trained on valid programs only, and a program's score is the negated
discriminator (or critic) output. This is a cheaper relative of AnoGAN~\cite{schlegl2017}. We compare it with Isolation
Forest, a one-class SVM, and the plain signal ``the compiler reported an error'', using AUROC~\cite{fawcett2006}.
